El \textit{World Editor} es una herramienta de desarrollo destinada a la creación y edición de mundos dentro del juego.
Está construido como un cliente de Unity que se comunica con el \textit{Core Service} a través del \textit{API Gateway}, siguiendo el mismo
patrón de acceso que el cliente de juego. El editor permite a los desarrolladores y diseñadores crear nuevos mundos,
modificar la configuración de zonas existentes, gestionar \textit{assets} y probar cambios en un entorno controlado antes de su despliegue en producción.

\subsubsection{Estructura de un Mundo}

Un mundo se compone de tres elementos principales.
\begin{itemize}
  \item \textbf{Zonas (\textit{Zones})}: Subdivisiones de un mundo, cada una alojada en una instancia de servidor independiente.
  Esta separación permite escalar horizontalmente y distribuir la carga entre múltiples servidores, mientras el estado
  del jugador se preserva de forma consistente al cambiar de zona.
  \item \textbf{Objetos Creables (\textit{Creatables})}: Objetos instanciados dinámicamente en tiempo de ejecución, como \textit{NPCs},
  enemigos o ítems. A diferencia de los \textit{Placeables}, no pueden ser posicionados manualmente durante la fase de diseño,
  sino que son generados por la lógica del juego.
  \item \textbf{Objetos Colocables (\textit{Placeables})}: Objetos posicionados manualmente en el editor durante la fase de diseño,
  que permanecen estáticos en tiempo de ejecución. Se clasifican en estructuras (\textit{Structures}), generadores
  (\textit{Spawners}) y objetos diversos (\textit{Miscellaneous Objects}).
\end{itemize}

\subsubsection{Zonas (\textit{Zones})}

En los juegos \textit{MMORPG} es habitual distribuir el mundo entre múltiples
servidores para poder sostener entornos de gran escala y una alta cantidad de jugadores simultáneos. Para abordar este
desafío, cada mundo se divide en zonas independientes, cada una ejecutándose en su propia instancia de servidor dentro
del clúster. El estado del jugador, incluyendo su inventario y progresión, se mantiene consistente a través de todas
las zonas de un mismo mundo, permitiendo que el jugador transite entre ellas sin perder información.
\vspace{0.25cm}

Los creadores de mundos pueden crear nuevas zonas dentro del editor, configurando parámetros como la textura
del terreno, la textura del cielo (\textit{Skybox}) y también configurar cómo un jugador se transporta de una zona a otra mediante
portales (\textit{Portals}).

\subsubsection{Objetos Colocables (\textit{Placeable Objects})}

Los \textit{Placeables} son objetos posicionados manualmente en el editor durante la fase de diseño y permanecen estáticos en
tiempo de ejecución. Se clasifican en tres categorías principales:

\begin{itemize}
  \item \textbf{Estructuras (\textit{Structures})}: Objetos estáticos que forman parte del entorno del juego, como edificios,
  árboles o rocas. Al ser colocados, su posición, rotación y escala pueden editarse mediante \textit{Gizmos} o a través
  de su panel de configuración. Cada estructura permite configurar sus colisionadores (\textit{Colliders}), que definen
  las interacciones físicas con el jugador y otros elementos del juego. Actualmente se soportan dos tipos de colisionador:
  cubo (\textit{Cube}), de uso general, y rampa (\textit{Slope}), un colisionador de forma piramidal diseñado para
  superficies inclinadas que permiten al jugador desplazarse hacia zonas de mayor altura. Adicionalmente, cualquier
  estructura puede configurarse como tienda (\textit{Shop}), habilitando la compra de ítems por parte de los jugadores.

  \item \textbf{Generadores (\textit{Spawners})}: Objetos que instancian entidades en tiempo de ejecución, como enemigos o \textit{NPCs}.
  Pueden configurarse con parámetros como frecuencia de aparición, rango y tipo de entidad a generar. Actualmente se
  encuentran disponibles tres tipos: generador de enemigos (\textit{Enemy Spawner}), generador de \textit{NPCs} (\textit{NPC Spawner})
  y generador de jugadores (\textit{Player Spawner}).

  \item \textbf{Objetos Diversos (\textit{Miscellaneous Objects})}: Elementos que no se ajustan a las categorías anteriores pero
  contribuyen a la composición del mundo. Actualmente se encuentran disponibles dos variantes:
  \begin{itemize}
    \item \textbf{Cofres (\textit{Chests})}: Objetos interactuables que los jugadores pueden abrir para obtener recompensas.
    Configurables con distintos tipos de botín, tiempo de reseteo y apariencia en estado abierto y cerrado.
    \item \textbf{Portales (\textit{Portals})}: Objetos que permiten transportar al jugador a otra ubicación, incluso entre zonas distintas
    del mundo. Al interactuar con un \textit{Portal}, el jugador es teletransportado a la zona de destino configurada,
    conservando su estado y progresión.
  \end{itemize}
\end{itemize}

\subsubsection{Objetos Creables (\textit{Creatable Objects})}

Los \textit{Creatables} son objetos instanciados dinámicamente en tiempo de ejecución. A diferencia de los \textit{Placeables}, no se
posicionan directamente en el mundo, sino que se configuran y gestionan a través de paneles de edición dedicados.
Actualmente se encuentran disponibles los siguientes tipos:

\begin{itemize}
  \item \textbf{\textit{NPC} Pasivo (\textit{Passive NPC})}: Personajes no hostiles con los que el jugador puede
  interactuar mediante diálogos, como aldeanos o comerciantes. El creador puede asignarles diálogos predefinidos
  para que el jugador obtenga información, inicie \textit{Quests} o simplemente enriquezca la ambientación del mundo.

  \item \textbf{\textit{NPC} Agresivo (\textit{Aggressive NPC})}: Personajes hostiles que atacan al jugador al
  detectarlo dentro de un rango determinado. Configurables con estadísticas de salud, tipos de armas que utilizan
  y tabla de botín que sueltan al ser derrotados.

  \item \textbf{Objetos Consumibles (\textit{Consumable Items})}: Ítems que el jugador puede recoger y utilizar durante
  el juego, como pociones de salud o comida. Se consumen al ser utilizados y pueden otorgar distintos efectos
  según su configuración.

  \item \textbf{Armas (\textit{Weapon Items})}: Ítems que el jugador puede recoger y equipar para atacar \textit{NPCs} agresivos
  u otros jugadores. Configurables con tipo de arma (cuerpo a cuerpo o a distancia), daño, alcance, velocidad
  de ataque y textura.

  \item \textbf{Tablas de Botín (\textit{Loot Tables})}: Configuraciones que definen qué ítems y qué cantidad de oro
  se obtienen al interactuar con un elemento del juego. Se asignan a enemigos y \textit{Chests}, donde cada ítem de la
  tabla tiene una probabilidad de aparición determinada.

  \item \textbf{Diálogos (\textit{Dialogs})}: Conjuntos de mensajes asignables a \textit{NPCs} pasivos que habilitan la
  interacción con el jugador. Permiten estructurar flujos de conversación y vincular \textit{Quests} a mensajes
  específicos, posibilitando mostrar distintos diálogos según el progreso del jugador.

  \item \textbf{Misiones (\textit{Quests})}: Misiones asignables a \textit{NPCs} que el jugador puede completar para obtener
  recompensas. Los objetivos disponibles actualmente incluyen eliminar una cantidad determinada de enemigos o
  interactuar con \textit{NPCs} específicos. Al completarse, el jugador puede recibir experiencia, oro o ítems especiales.

  \item \textbf{Tienda (\textit{Shop})}: Configuración asignable a estructuras que habilita la compra de ítems al
  interactuar con ellas. Soporta dos tipos de productos: ítems consumibles y armas, adquiribles con oro
  obtenido durante el juego, y cosméticos, adquiribles con gemas.

  \item \textbf{Cosméticos (\textit{Cosmetics})}: Objetos de personalización estética como ropa, accesorios o \textit{skins}
  que no afectan la jugabilidad. Se adquieren con gemas, una moneda \textit{premium} obtenida con dinero real, y
  pueden configurarse en tiendas con distintos precios y cantidades disponibles.
\end{itemize}

\subsubsection{Gestión de Mundos}

El editor cuenta con un sistema de guardado y apertura de mundos similar al de otros programas de edición,
permitiendo retomar la edición en distintas sesiones de trabajo.

Los mundos se persisten localmente utilizando \texttt{Application.persistentDataPath} de Unity, una ruta
del sistema de archivos gestionada por el motor que garantiza permisos de escritura en todas las plataformas
soportadas. Cada mundo se almacena como un directorio independiente con el nombre del mundo, con la
siguiente estructura:

\begin{verbatim}
<mundo>/
├── world.data
├── creatables.data
└── zone_1.world
    zone_2.world
    ...
    zone_n.world
\end{verbatim}

\begin{itemize}
  \item \texttt{world.data}: Contiene la configuración general del mundo, incluyendo nombre, descripción
  e identificador asignado al momento de su publicación.
  \item \texttt{creatables.data}: Almacena la definición de todos los objetos creables configurados
  para el mundo, como \textit{NPCs}, enemigos, ítems y misiones.
  \item \texttt{zone\_\{n\}.world}: Representa el estado de una zona independiente, almacenando todos
  los \textit{Placeables} colocados en ella junto con su posición, rotación, escala y propiedades configuradas.
  Se genera un archivo por cada zona que compone el mundo.
\end{itemize}

\subsubsection{Personalización del Mundo}

El editor permite a los creadores personalizar la apariencia de cada zona mediante la importación de contenido
propio, incluyendo modelos 3D, texturas de ítems, texturas de suelo, \textit{skyboxes} y cosméticos personalizados
para el jugador.
\vspace{0.25cm}

Para quienes no deseen importar contenido propio, el editor ofrece la posibilidad de descargar un conjunto
de \textit{assets} predeterminados que incluye modelos 3D para estructuras, enemigos, \textit{NPCs}, ítems y cosméticos,
así como texturas de suelo y cielo para las zonas.
\vspace{0.25cm}

La importación de modelos 3D en tiempo de ejecución, funcionalidad no soportada nativamente por Unity,
se resolvió mediante la integración de la librería \texttt{GLTFast} % TODO(reference)
, que permite cargar modelos en formato \texttt{.glb} durante la ejecución del editor, facilitando
la incorporación de \textit{assets} personalizados por parte de los creadores de mundos.

\subsubsection{Economía}

El sistema cuenta con paneles de gestión económica destinados a los creadores de mundos, que centralizan
el seguimiento de sus mundos publicados. Desde esta sección, el creador puede visualizar el
estado de cada mundo publicado, ver la fecha de renovación de su período de suscripción y
administrar sus publicaciones.
\vspace{0.25cm}

Adicionalmente, existe un panel de ingresos que permite a los creadores recaudar las
ganancias generadas en sus mundos. Dichos ingresos provienen de las compras realizadas por jugadores
con gemas dentro del mundo, correspondientes a la adquisición de cosméticos disponibles en las tiendas
configuradas por el creador. Más información sobre los flujos de zonas se encuentra explicados en la sección 
\nameref{sec:payment-service}, mientras que los detalles de cómo funciona la economía se encuentran en la sección 
\nameref{sec:business-model}.

\subsubsection{Gestión de Recursos en Memoria}

La gestión de recursos en memoria se organiza mediante un sistema de \textit{managers} y repositorios internos.
El \texttt{DataPersistenceManager} centraliza las operaciones de guardado y carga, delegando la
persistencia de cada tipo de dato a un repositorio especializado: \texttt{ZonesRepository},
\texttt{WorldRepository} y \texttt{CreatablesRepository}, entre otros. Cada repositorio encapsula
la lógica de serialización y acceso a disco de su dominio correspondiente.
\vspace{0.25cm}

Para los \textit{Placeables}, cada objeto colocado en el mundo existe como una instancia rastreada en la jerarquía
de la escena de Unity. Mediante un sistema de eventos, cada instancia notifica al \textit{manager} correspondiente
cuando su estado ha sido modificado, determinando así si debe ser incluida en el próximo ciclo de guardado.
\vspace{0.25cm}

Los \textit{Creatables}, al no tener representación física en la escena, son gestionados por un
\texttt{CreatablesManager} dedicado que mantiene en memoria el conjunto de objetos creables del mundo,
permitiendo su consulta, modificación y persistencia de forma centralizada.
\vspace{0.25cm}

Cada \textit{Placeable} del mundo es un objeto instanciado y rastreado en la escena del juego, que hereda de la
clase abstracta genérica \texttt{Placeable<T>}, donde \texttt{T} representa el modelo de datos asociado
al objeto, definido en el \textit{Shared Package}. Esta clase base implementa dos interfaces clave:
\texttt{IPersistent<ZoneData>}, que obliga a cada \textit{Placeable} a definir cómo serializa su estado en los
datos de la zona mediante \texttt{SaveData}, e \texttt{ILoadable<T>}, que define cómo reconstruye su
estado a partir de los datos cargados mediante \texttt{LoadData}.

Al ser instanciado en la escena, cada \textit{Placeable} se registra en el sistema mediante el método
\texttt{RegisterPlaceable}, que emite un evento (\texttt{ZoneDataRegistryEvent}) capturado por el
\texttt{DataPersistenceManager}. De esta forma, el \textit{manager} mantiene un registro actualizado de todos
los objetos presentes en la zona sin necesidad de búsquedas en la jerarquía de la escena, delegando
la lógica de serialización a cada objeto.

Como ejemplo concreto, \texttt{StructureObject} extiende \texttt{Placeable<StructureData>} y define
en \texttt{SaveData} cómo persiste su posición, rotación, escala y configuración de colisionadores
en el \texttt{ZoneData} de la zona activa, mientras que \texttt{LoadData} reconstruye el objeto en
la escena a partir de los datos almacenados.

\subsubsection{Publicación de Mundos}

Una vez finalizado el diseño, el creador puede publicar su mundo para que pueda ser accedido mediante el cliente de juego.
El proceso requiere haber iniciado sesión y se realiza desde el panel de publicación. Al confirmar la publicación,
se inicia una secuencia de carga hacia el \textit{Core Service} que sigue el siguiente orden:

\begin{enumerate}
  \item \textbf{Datos del Mundo (\textit{World Data})}: Los datos estructurales del mundo en formato \textit{JSON}, que incluyen la configuración
  de zonas, objetos y sus propiedades.
  \item \textbf{Modelos 3D}: Los modelos utilizados en el mundo, excluyendo aquellos que forman parte del
  conjunto de \textit{assets} predeterminados del sistema.
  \item \textbf{\textit{Assets} de Ítems}: Las imágenes e iconos asociados a los ítems creados por el usuario.
  \item \textbf{Cosméticos}: Los \textit{assets} visuales correspondientes a los cosméticos configurados en el mundo.
  \item \textbf{\textit{Assets} de Entorno}: Las texturas de suelo (\textit{ground textures}) y cielos
  (\textit{skybox textures}) asociadas a cada zona del mundo.
\end{enumerate}